% Supplement to p2.tex. Its theorem, section and equation numbers are those of p2.pdf, read from p2.aux;
% build p2.tex first (make papers does).
\newcommand{\papertitle}{Supplement to ``Increasing ratios in finite nonnegative linearization'': formalization and AI use}
\newcommand{\papershort}{SUPPLEMENT: INCREASING RATIOS IN FINITE LINEARIZATION}
\newcommand{\paperdate}{November 1, 2026}
\documentclass[twoside,10pt]{article}
\usepackage[T1]{fontenc}
\usepackage{lmodern}
\usepackage{amsmath,amssymb,amsthm,amsfonts,mathtools}
\usepackage[letterpaper,total={5.5in,8in},vmarginratio=1:1,hmarginratio=1:1,headsep=21pt]{geometry}
\usepackage{xcolor}
\usepackage{array,longtable}
\usepackage{xr-hyper}
\usepackage[breaklinks=true,linktocpage=true]{hyperref}
\hypersetup{colorlinks=true,linkcolor=blue!45!black,urlcolor=blue!45!black,citecolor=blue!45!black,
pdfauthor={Leslie P. Polzer},pdftitle={\papertitle},
pdfsubject={Finite orthogonal polynomials and nonnegative linearization},
pdfkeywords={orthogonal polynomials, Krawtchouk polynomials, nonnegative linearization, quadratic transformation, hypergroups, Lean}}
\medmuskip=4mu
\thickmuskip=8mu
\setlength{\emergencystretch}{2em}
\numberwithin{equation}{section}
\newtheoremstyle{slbody}{8pt}{8pt}{\normalfont\slshape}{}{\bfseries}{.}{0.5em}
{\thmname{#1}\thmnumber{ #2}\thmnote{ {\normalfont\itshape(\ignorespaces#3\unskip)}}}
\theoremstyle{slbody}
\newtheorem{theorem}{Theorem}[section]
\newtheorem{lemma}[theorem]{Lemma}
\newtheorem{corollary}[theorem]{Corollary}
\newtheorem{proposition}[theorem]{Proposition}
\theoremstyle{definition}
\newtheorem{definition}[theorem]{Definition}
\theoremstyle{remark}
\newtheorem{remark}[theorem]{Remark}
\renewcommand{\qedsymbol}{$\blacksquare$}
\newcommand{\qint}[1]{[#1]_q}
% \xnum{P1-thm:lc}{3.1} prints the fixed number of a result in a companion paper; check_xref.py checks it.
\newcommand{\xnum}[2]{#2}
\title{\papertitle}
\author{Leslie P. Polzer}
\newcommand{\authoraffiliation}{Independent Researcher}
\newcommand{\authoremail}{polzer@fastmail.com}
\date{\paperdate}
\makeatletter
\renewcommand{\section}{\@startsection{section}{1}{\z@}{-3.5ex \@plus -1ex \@minus -.2ex}{2.3ex \@plus .2ex}{\normalfont\large\bfseries}}
\renewcommand{\@seccntformat}[1]{\csname the#1\endcsname.\enspace}
\def\tagform@#1{\maketag@@@{(\ignorespaces{\oldstylenums{#1}}\unskip\@@italiccorr)}}
\renewcommand{\eqref}[1]{\textup{{\normalfont(\oldstylenums{\ref{#1}}}\normalfont)}}
\newcommand{\shorttitle}{\papershort}
\newcommand{\shortauthors}{LESLIE P. POLZER}
\newcommand{\ps@bookheader}{%
\renewcommand\@oddfoot{\hfil}\renewcommand\@evenfoot{\hfil}%
\renewcommand\@oddhead{\ifnum\value{page}>1{\small\hfil\shortauthors}\hfil\thepage\else\hfil\fi}%
\renewcommand\@evenhead{\thepage\hfil\small\shorttitle\hfil}}
\pagestyle{bookheader}
\renewcommand{\maketitle}{\begin{center}{\large\bfseries\@title}\\\vskip36pt
{\scshape\@author}\par\medskip{\small\authoraffiliation}\par\smallskip
{\small\texttt{\authoremail}}\par\smallskip{\small\@date}\par\vskip36pt\end{center}}
\makeatother
\renewenvironment{abstract}{\quotation\noindent\small{\bfseries\abstractname.\enspace}}{\endquotation}

\externaldocument[M-]{p2}
\begin{document}
\maketitle

\noindent This supplement to~\cite{main} gives the Lean declaration for each of its results, the inventory of its computer-assisted steps, the statements that are not formalized, and the full AI-use disclosure. Theorem, section, equation and appendix numbers refer to~\cite{main}.

\section*{Code and formalization}
The companion code~\cite{code} contains the certificate scripts cited in the proofs of Theorems~\ref{M-thm:lc-general} and~\ref{M-thm:lc-cubic}, exact-arithmetic diagnostics for every result, and a Lean~4 development. Its axiom audit rejects dependencies outside \path{propext}, \path{Classical.choice} and \path{Quot.sound}; the declarations below have compiled proofs without added axioms or admitted goals. The documentation revision~\cite{docs} describes the diagnostics, explains how each Lean statement matches the text, and records a run of all scripts from a fresh clone.
\begin{longtable}{@{}>{\raggedright\arraybackslash}p{0.36\linewidth}>{\raggedright\arraybackslash}p{0.59\linewidth}@{}}
Result & Lean declaration\\[3pt]
\endhead
Theorem~\ref{M-thm:lc-vary}: varying ratio & \path{Szkraw.varyratio_matrix}\\[3pt]
Corollary~\ref{M-cor:lc-vary}: nonincreasing weights & \path{Szkraw.varyratio_weights}\\[3pt]
Proposition~\ref{M-prop:lc-quad}: quadratic arrays & \path{Szkraw.quadratic_matrix}\\[3pt]
Theorem~\ref{M-thm:lc-affine}: linearly increasing ratio & \path{Szkraw.affine_matrix}, \path{Szkraw.affine_bmono_iff}\\[3pt]
Remarks after Theorem~\ref{M-thm:lc-affine} & \path{Szkraw.affine_range_remarks}, \path{Szkraw.profile_last_le}, \path{Szkraw.affEx_remark}, \path{Szkraw.affine_linear_quadratic}\\[3pt]
Theorem~\ref{M-thm:lc-quadratic}: convex quadratic ratio & \path{Szkraw.quadratic_ratio_matrix}, \path{Szkraw.quadratic_bmono_iff}\\[3pt]
Remarks on Theorem~\ref{M-thm:lc-quadratic} & \path{Szkraw.quadExample_remark}, \path{Szkraw.quadB_unit_fails}\\[3pt]
Theorem~\ref{M-thm:lc-lowdeg}: low-degree profiles & \path{Szkraw.lowdeg_matrix}\\[3pt]
Corollary~\ref{M-cor:lc-lowdeg} & \path{Szkraw.lowdeg_newton_lc}, \path{Szkraw.lowdeg_ratio_matrix}\\[3pt]
Proposition~\ref{M-prop:np}: products of Newton sequences & \path{Szkraw.newton_product_LCNZ}\\[3pt]
Lemma~\ref{M-lem:abel-weights} & \path{Szkraw.abel_weight}\\[3pt]
Theorem~\ref{M-thm:lc-general}: log-concave ratio & \path{Szkraw.lc_general}\\[3pt]
Theorem~\ref{M-thm:lc-cubic}: cubic ratio & \path{Szkraw.cubic_ratio_matrix}\\[3pt]
Corollary~\ref{M-cor:cubic74}: extension to $r\sigma\leq\frac74t\tau$ & \path{Szkraw.cubic_ratio74_matrix}\\[3pt]
Appendix~\ref{M-app:frontier}: the counterexamples with $N=3$ and the cubic counterexample with $N=4$ & \path{Szkraw.openEx1_fails}, \path{Szkraw.openEx2_fails}, \path{Szkraw.cubic_ratio_fails}\\[3pt]
Appendix~\ref{M-app:frontier}: the counterexamples with $N=48$, $64$, $40$, $40$ and $120$ & \path{Szkraw.cubicBinom48_fails}, \path{Szkraw.cubicPow64_fails}, \path{Szkraw.steepRatio40_fails}, \path{Szkraw.topRatio40_fails}, \path{Szkraw.topRatio120_fails}\\[3pt]
Remarks after Corollary~\ref{M-cor:lc-vary} & \path{Szkraw.strict_ratio_example}, \path{Szkraw.openEx1_b_newton}, \path{Szkraw.linear_ratio_concave_b}\\[3pt]
After Proposition~\ref{M-prop:lc-quad}: each hypothesis is needed & \path{Szkraw.quadratic_hypotheses_needed}\\[3pt]
Section~\ref{M-sec:higher}: exponential ratios, the ratio of the cubic counterexample, containments and non-comparability & \path{Szkraw.exponential_ratio}, \path{Szkraw.cubic_example_ratio_newton}, \path{Szkraw.general_contains}, \path{Szkraw.quad_lowdeg_incomparable}\\[3pt]
After Lemma~\ref{M-lem:abel-weights}: weight sums on geometric sequences & \path{Szkraw.weight_sums_geometric}\\[3pt]
Proof of Theorem~\ref{M-thm:lc-general}: $(\log\rho)''(0)>0$ for $\rho(x)=1+x/2+x^2/2$ & \path{Szkraw.log_rho_not_concave}\\[3pt]
Remark~\ref{M-rem:cubic-limit}: geometric profiles and the case $u>2$ & \path{Szkraw.geometric_kappa_values}, \path{Szkraw.geometric_kappa}\\[3pt]
Appendix~\ref{M-app:frontier}: a negative entry of the operator $M$ for $N=6$, $t=\tau=0$ & \path{Szkraw.opM_negative_entry}, \path{Szkraw.opM_has_negative_entry}\\[3pt]
After Proposition~\ref{M-prop:np}: the law of $|S\cup T|$, ultra-log-concavity of $|S\cup T|$, and the double-binomial convolution $h_k=\kappa_k$ & \path{Szkraw.RandomSets.union_card_law}, \path{Szkraw.RandomSets.random_sets_ulc}, \path{Szkraw.RandomSets.walkup_convolution}
\end{longtable}
Not formalized: the identification of the random-set reading after Proposition~\ref{M-prop:np} with Walkup's formulation (his conditionally uniform random sets and their fundamental sequences), which concerns the literature; the reported failure of \path{runs/cubic/bdetour_u.py} at $u=9/5$, which concerns one certificate search, not property (P); and, in Appendix~\ref{M-app:frontier}, the sampled tests of cubic ratios and of the resummed operator, which are numerical evidence over finitely many arrays. More generally, attributions to other works and statements about what they contain are not formalized, for instance that Proposition~\ref{M-prop:np} is Walkup's Theorem~2 in the Newton basis and that the Abel argument is that of Wang and Yeh: Lean checks our statements, not the cited papers.

\section*{Certificate inventory}
Table~\ref{tab:certs} lists every computer-assisted step of Theorems~\ref{M-thm:lc-general} and~\ref{M-thm:lc-cubic} and of Corollary~\ref{M-cor:cubic74}. ``Symbolic'' means a certificate with polynomial or rational coefficients that covers the whole stated range; ``finite'' means one exact certificate for each listed case. Scripts are in~\cite{code}, under \path{runs/lean_plan/} or \path{runs/cubic/}; Lean names omit the prefix \path{Szkraw.}; the documentation~\cite{docs} describes each script.

\begin{longtable}{@{}>{\raggedright\arraybackslash}p{0.22\linewidth}>{\raggedright\arraybackslash}p{0.24\linewidth}>{\raggedright\arraybackslash}p{0.19\linewidth}>{\raggedright\arraybackslash}p{0.25\linewidth}@{}}
\caption{Computer-assisted steps.}\label{tab:certs}\\
Step & Range & Script & Lean declaration\\\hline
\endfirsthead
Step & Range & Script & Lean declaration\\\hline
\endhead
\multicolumn4{@{}l}{\emph{Theorem~\ref{M-thm:lc-quadratic}} (no computer assistance; for comparison)}\\
\eqref{M-eq:quad-cert} & $n\geq2$; printed table for $n\leq6$ & none & \path{cert_nonneg}\\[3pt]
\multicolumn4{@{}l}{\emph{Theorem~\ref{M-thm:lc-general}} (\path{lc_general})}\\
(vi): $e_n\geq0$ for $K$, $L$, $R$, $R_2$, $K''$, $Q^*$, $K_a$, $\Psi$, $\widehat Q$ & all $n$, all $a\geq3$ ($a\geq5$ for $K_a$, $a\geq6$ for $R_2$); symbolic & \path{edge_cert.py}; exact check for $a\leq60$: \path{edge_exact.py} & \path{eK_nonneg}, \path{eL1_nonneg}, \path{eR_nonneg}, \path{eR2_nonneg}, \path{eKpp_nonneg}, \path{eQs_nonneg}, \path{eKa_nonneg}, \path{ePsi_nonneg}, \path{eQhat_nonneg}\\[3pt]
(vii): $M_{1,m}+M_{m,1}\geq0$; $S_{2,0}\geq0$ & $3\leq a\leq5$, $2\leq m\leq a$; $a\in\{3,4\}$; finite & \path{window_cert.py} & \path{win_MM1_3_2}, \dots, \path{win_MM1_5_5}; \path{win_S20_3}, \path{win_S20_4}\\[3pt]
\multicolumn4{@{}l}{\emph{Theorem~\ref{M-thm:lc-cubic} and Corollary~\ref{M-cor:cubic74}} (\path{cubic_ratio74_matrix}; \path{cubic_ratio_matrix} follows)}\\
\eqref{M-eq:cubic-B} & $n\geq6$, all $\ell$; $3\leq n\leq5$, all $\ell$; $n=2$, $\ell\geq1$; symbolic in $\ell$ and $n$ & \path{gen_cubicB.py}; for $n\leq5$ also \path{cert74.py} & \path{cubicB_gen}; \path{cubicB_three}, \path{cubicB_four}, \path{cubicB_five}; \path{cubicB_two}\\[3pt]
\eqref{M-eq:cubic-corner} & $\ell=0$, $n=2$ & \path{gen_cubicB.py} & \path{cubicB_corner}\\[3pt]
\eqref{M-eq:cubic-E} & $n\geq7$, all $\ell$; $3\leq n\leq6$, all $\ell$; symbolic & \path{gen_cubicE.py} & \path{cubicE_gen}; \path{cubicE_three}, \dots, \path{cubicE_six}\\[3pt]
$\Phi^{\mathrm{loc}}_3\geq0$ in \eqref{M-eq:cubic-key} & all $n$, all $a\geq3$; symbolic & \path{phi3_generic.py}; exact check: \path{phi3_edge.py} & \path{Phi3_nonneg}\\[3pt]
Corollary~\ref{M-cor:cubic74}, written route: $\kappa_{t\tau}+\frac74\kappa_{r\sigma}$ in \eqref{M-eq:cubic-B}, \eqref{M-eq:cubic-E} & $n\leq5$; regimes $z\in\{0,\ell,1/q\}$; symbolic & \path{bdetour_u.py}, \path{edetour_u.py} & Lean uses the route of \path{cert74.py} (rows above)
\end{longtable}

\section*{Verification procedure}
The certificates of Table~\ref{tab:certs} can be checked in two independent ways from a clone of the release~\cite{code}.

\emph{Lean.} From the repository root, \texttt{make lean} builds the development and \texttt{make axioms} writes the axiom audit \path{lean/axioms.log} of the principal statements, among them \path{Szkraw.lc_general} and \path{Szkraw.cubic_ratio74_matrix}. The audit of a statement covers every declaration it depends on, and so every declaration of Table~\ref{tab:certs}; it must report only \path{propext}, \path{Classical.choice} and \path{Quot.sound}. A declaration that compiles proves its obligation for every $N$; Lean does not rely on the output of any script.

\emph{Scripts.} The scripts need Python~3 with SymPy, NumPy and SciPy, and are run from their own directory, \path{runs/lean_plan/} unless stated otherwise, for example \texttt{python3 window\_cert.py}. The commands are:
\begin{center}
\renewcommand{\arraystretch}{1.15}
\begin{tabular}{@{}>{\raggedright\arraybackslash}p{0.3\linewidth}>{\raggedright\arraybackslash}p{0.64\linewidth}@{}}
Obligation & Commands\\\hline
step (vi) & \texttt{edge\_exact.py 60} (exact cross-check, $a\leq60$); \texttt{edge\_cert.py} (symbolic numerators)\\
step (vii) & \texttt{window\_cert.py}\\
\eqref{M-eq:cubic-B}, \eqref{M-eq:cubic-corner} & \texttt{gen\_cubicB.py C OUT} for \texttt{C} in \texttt{gen}, \texttt{2}, \texttt{3}, \texttt{4}, \texttt{5}, \texttt{corner}\\
\eqref{M-eq:cubic-E} & \texttt{gen\_cubicE.py C OUT} for \texttt{C} in \texttt{gen}, \texttt{3}, \texttt{4}, \texttt{5}, \texttt{6}\\
$\Phi^{\mathrm{loc}}_3\geq0$ & \texttt{phi3\_generic.py}; \texttt{phi3\_edge.py 40} (exact cross-check)\\
Corollary~\ref{M-cor:cubic74}, written route & in \path{runs/cubic/}: \texttt{bdetour\_u.py 7/4}, \texttt{edetour\_u.py 3 7/4}
\end{tabular}
\end{center}
Each script prints its verdict, for instance \texttt{ALL CERTIFIED} from \texttt{window\_cert.py} and an empty failure list for each form from \texttt{edge\_exact.py}; the exit code alone is not a verdict. The scripts \texttt{gen\_cubicB.py} and \texttt{gen\_cubicE.py} also write the Lean certificate for the case \texttt{C} to \texttt{OUT}, which can be compared with the corresponding file in \path{lean/Szkraw/}. The record~\cite{docs} lists a run of all these commands except \texttt{edge\_cert.py} from a fresh clone of revision \texttt{2f96f4d}, with the hash of each script and its last output line.

\section*{AI-use disclosure}
The following AI models were used: Claude Opus 5.5 and, for parts of the Lean formalization of Theorem~\ref{M-thm:lc-general}, Claude Haiku (Anthropic; the Haiku version was not recorded), both run in the Claude Code agent; GPT-6 Astra or GPT-6.1 Sol (OpenAI; which of the two was not recorded), run in the Codex agent; and GPT-6 Pro (OpenAI), used in ChatGPT. Claude Opus 5.5 and the Codex model were used to develop proof arguments, check source passages, construct exact-arithmetic diagnostics, draft and revise the paper, and prepare the Lean formalization. The resummation in the proof of Theorem~\ref{M-thm:lc-affine}, its key inequality under nondecreasing $b$, and the weighted detour and key inequality in the proof of Theorem~\ref{M-thm:lc-quadratic} arose in research consultations with GPT-6 Pro; the exact checks were completed with Claude Opus 5.5, which also prepared the Lean formalizations. Theorem~\ref{M-thm:lc-lowdeg}, Corollary~\ref{M-cor:lc-lowdeg}, Theorems~\ref{M-thm:lc-general} and~\ref{M-thm:lc-cubic} and Lemma~\ref{M-lem:abel-weights} were found and proved with Claude Opus 5.5, without research consultation; Proposition~\ref{M-prop:np}, which is due to Walkup, was rederived with Claude Opus 5.5 together with the self-contained proof given here. GPT-6 Pro reviewed every result, as part of a longer manuscript and with independent exact checks, and the corrections it requested were applied. For Theorems~\ref{M-thm:lc-general} and~\ref{M-thm:lc-cubic} the reviewer reran transcribed copies of the certificate scripts; the scripts were also run unmodified from a fresh clone of the companion code at the cited revision, and all of them passed~\cite{docs}. This AI review is not human peer review, and no independent human verification is claimed. The author set the research objective and is responsible for publication decisions and the final mathematical claims. ``Machine-verified'' refers only to the Lean theorems listed in the section on code and formalization; Theorems~\ref{M-thm:lc-general} and~\ref{M-thm:lc-cubic} and Corollary~\ref{M-cor:cubic74} are machine-verified.

\begin{thebibliography}{99}

\bibitem{main} L. P. Polzer, \emph{Increasing ratios in finite nonnegative linearization}, 2026, in release~1.0.6 of the companion code~\cite{code}, \url{https://github.com/skypher/szkraw/releases/tag/v1.0.6}.
\bibitem{code} L. P. Polzer, \emph{Companion code for finite nonnegative linearization}, release~1.0.6 (2026), \url{https://github.com/skypher/szkraw/releases/tag/v1.0.6}. This release is the single landing point for the evidence: it contains the three papers and their supplements, the Lean development and the scripts, which are unchanged since revision \texttt{2f96f4d}, \url{https://github.com/skypher/szkraw/tree/2f96f4d228215b35e31729f1d1526cd31879f161}, and the documentation of~\cite{docs}.
\bibitem{docs} L. P. Polzer, documentation of the companion code~\cite{code}: \path{papers/FORMALIZATION.md}, \path{papers/DIAGNOSTICS.md} and the record \path{runs/verify/RESULTS.md} of a run of all scripts from a fresh clone of revision \texttt{2f96f4d}, as of revision \texttt{4b2d7e1}, \url{https://github.com/skypher/szkraw/tree/4b2d7e13d5c9a9fe8bbd406959d0ecd2cc192de8}.
\end{thebibliography}
\end{document}
